%% toc: P15 163 163.1,163.2,163.3,163.4
% Worked boundaries for the definitions, canonical-reduction lemma, symmetrization
% theorem, directing-law space and finite-urn bound of the exchangeability chapter.
\chapter{Boundaries of the Exchangeable Fields, Symmetrized Averages, Directing Laws and Finite-Urn Bound}
\label{sec:p15-w02c-p15-022}

This chapter works out, for the entries of
\cref{sec:p15-exchangeability-and-definetti}, the cases in which each
definition or hypothesis reaches its edge. The objects are those of
\cref{def:p15-exchangeable-canonical}: the invariant fields
\(\mathcal E_n\), their intersection \(\mathcal E\), the tail field
\(\mathcal T\), exchangeability of a law, and the averages
\(A_{n,k}(\varphi)\) over ordered tuples of distinct indices; and those of
\cref{def:p15-directing-law}: the evaluation sigma-field on
\(\mathscr P(S)\), the countable determining class, the random law
\(K\), its information \(\sigma(K)\) and the null-augmented field
\(\mathcal G^P\); and the total variation bound
\(k(k-1)/(2N)\) of \cref{thm:p15-finite-urn-approximation} for draws
without replacement from a finite binary urn.

The \emph{inputs} throughout are a measurable value space
\((S,\mathcal S)\), the canonical sequence space with coordinates
\(X_j\), and a law \(P\) on it. The \emph{conclusions} are of three kinds:
an inclusion stated in the chapter is strict (or collapses), a
hypothesis of a result is shown to be needed by a law or a space for
which the conclusion fails, or a conclusion is shown to survive in a
weakened, explicitly stated form when the hypothesis is dropped.
Each result below carries a complete proof in the annex.

Durrett defines \(\mathcal E_n\) as the sigma-field generated by the
events invariant under the permutations that fix every index beyond
\(n\), and \(\mathcal E\) as their intersection
\cite[p.~250]{DurrettPTEAuthor2019}. Since the invariant events already
form a sigma-field (\cref{thm:p15-exchangeable-canonical}), that
generated sigma-field is exactly the class \(\mathcal E_n\) of
\cref{def:p15-exchangeable-canonical}; the two definitions name the same
object.

\section{The invariant fields at their edges}
\label{sec:p15-w02c-p15-022-fields}

The inclusions \(\mathcal E_{n+1}\subseteq\mathcal E_n\) and
\(\mathcal T\subseteq\mathcal E\) of
\cref{thm:p15-exchangeable-canonical} are exact inclusions of
collections of events. The first theorem decides when they are strict,
both literally and after adjoining null events, and shows which
hypothesis makes the tail and invariant information agree modulo null
events.

\begin{hbtheorem}{thm:p15-w02c-p15-022-field-strictness}{proved=annex:p15-w02c-p15-022-field-strictness}{When the invariant inclusions are strict}
Work in the setting of \cref{def:p15-exchangeable-canonical}.
\begin{enumerate}
\item If \(\mathcal S=\{\varnothing,S\}\), then
      \(\mathcal F=\{\varnothing,\Omega\}\), and
      \(\mathcal E_n=\mathcal E=\mathcal T=\mathcal F\) for every
      \(n\geq1\). No inclusion is strict.
\item If some \(B\in\mathcal S\) has \(\varnothing\neq B\neq S\), put
      \[
       A_n=\{X_j\in B\text{ for }1\leq j\leq n\},
      \]
      and let \(A\) be the event that exactly one index \(j\geq1\) has
      \(X_j\in B\).
      Then \(A_n\in\mathcal E_n\setminus\mathcal E_{n+1}\) for every
      \(n\geq1\), and \(A\in\mathcal E\setminus\mathcal T\). Every
      inclusion \(\mathcal E_{n+1}\subseteq\mathcal E_n\) and the
      inclusion \(\mathcal T\subseteq\mathcal E\) are strict.
\item Let \(P\) be exchangeable and \(B,A_n\) as in item~2. If
      \(P(X_1\in B,\ldots,X_n\in B,\ X_{n+1}\notin B)>0\), then
      \(A_n\notin\mathcal E_{n+1}^{P}\); the strictness of item~2
      survives modulo null events. This holds for every i.i.d.\ law
      whose coordinate law \(\mu\) has \(0<\mu(B)<1\). If instead \(P\)
      is the law of \((\xi,\xi,\ldots)\) for one \(S\)-valued random
      element \(\xi\), then \(A_n\in\mathcal E_{n+1}^{P}\) for every
      \(n\): the literal strictness disappears modulo null events.
\item Let \(S=\{0,1\}\) with all subsets measurable, and let \(P\) give
      probability \(1/2\) to each of the sequences
      \(a=(1,0,0,\ldots)\) and \(b=(0,0,0,\ldots)\). This law is not
      exchangeable, every event of \(\mathcal T\) has probability zero
      or one, and the event \(A\) of item~2 with \(B=\{1\}\) has
      \(P(A)=1/2\). Hence \(A\in\mathcal E\setminus\mathcal T^{P}\) and
      \(\mathcal E^{P}\neq\mathcal T^{P}\).
\end{enumerate}
\end{hbtheorem}

Item~4 shows that the equality \(\mathcal E^{P}=\mathcal T^{P}\) of
\cref{thm:p15-directing-law-uniqueness} needs its exchangeability
hypothesis: without it, invariant information can exceed tail
information by an event of probability \(1/2\). The law of item~4 is the
one used for the Hewitt--Savage countermodel in item~4 of
\cref{thm:p15-zero-one-countermodels}; here it is read as a boundary of
the invariant fields themselves. Items~2 and~3 together separate two
questions that the inclusions leave open: whether an invariant event of
level \(n\) can fail to be invariant at level \(n+1\) (it can, whenever
\(\mathcal S\) is nontrivial), and whether a law can see that failure
(it does exactly when the displayed configuration has positive
probability).

\begin{hbtheorem}{thm:p15-w02c-p15-022-not-exchangeable}{proved=annex:p15-w02c-p15-022-partial}{Partial symmetry and the canonical reduction}
Let \(S=\{0,1\}\) with all subsets measurable.
\begin{enumerate}
\item Let \(\xi\) be a fair bit, and let \(X_{2m-1}=\xi\) and
      \(X_{2m}=1-\xi\) for \(m\geq1\). Every \(X_j\) is a fair bit, and
      \((X_1,X_2)\) has the same law as \((X_2,X_1)\). The vector
      \((X_1,X_2,X_3)\) does not have the same law as
      \((X_1,X_3,X_2)\), so the law \(P\) of the sequence is not
      exchangeable. For the transposition \(\tau\) of \(1\) and \(2\),
      \[
       P(X_1=X_3)=1,\qquad
       P\bigl(T_\tau^{-1}\{X_1=X_3\}\bigr)=P(X_2=X_3)=0,
      \]
      so \(P\circ T_\tau^{-1}\neq P\) although the vector \((X_1,X_2)\)
      is symmetric under \(\tau\).
\item Let \(P\) be the law of \((\xi,\xi,\ldots)\) for a fair bit
      \(\xi\). On \(\Omega'=\Omega\times\{0,1\}\) with product
      sigma-field put \(Z(\omega,u)=\omega\) and \(U(\omega,u)=u\). Let
      \(Q_1=P\otimes\operatorname{Bernoulli}(1/2)\), and let \(Q_2\) be
      the image of \(P\) under \(\omega\mapsto(\omega,\omega_1)\). Both
      laws give \(Z\) the law \(P\). With
      \(\bar Y=\limsup_nn^{-1}\sum_{j=1}^nX_j\),
      \[
       E_{Q_1}\bigl(U\mid Z^{-1}(\mathcal E)\bigr)=\tfrac12,
       \qquad
       E_{Q_2}\bigl(U\mid Z^{-1}(\mathcal E)\bigr)=\bar Y\circ Z,
      \]
      almost surely under the respective laws, and \(\bar Y\circ Z\)
      takes the values \(0\) and \(1\) with probability \(1/2\) each
      under \(Q_2\).
\end{enumerate}
\end{hbtheorem}

Item~1 is the boundary of the equivalence in
\cref{thm:p15-exchangeable-canonical}. Exchangeability quantifies over
every finite coordinate vector. Symmetry of one initial vector under the
permutations of its own indices does not make those permutations
preserve the law of the whole sequence, because a permutation of
\(\{1,2\}\) also changes the joint law of the moved coordinates with the
later ones. The proof of the lemma uses vectors longer than the support
of the permutation for exactly this reason.

Item~2 is the boundary of the reduction in the same lemma. The reduction
transfers every identity that concerns the sequence and its pulled-back
invariant information, and nothing more: two probability spaces carrying
sequences with the same law \(P\) can give different conditional
expectations, given that pulled-back information, of a variable that is
not a function of the sequence.

\section{The symmetrized averages at their edges}
\label{sec:p15-w02c-p15-022-averages}

\begin{hbtheorem}{thm:p15-w02c-p15-022-average-edges}{proved=annex:p15-w02c-p15-022-average-edges}{The index range and the distinctness of indices}
Let \(P\) be exchangeable.
\begin{enumerate}
\item \emph{The index range.} For \(n<k\) the set \(I_{n,k}\) is empty
      and \((n)_k=0\), so \(A_{n,k}(\varphi)\) is undefined. At \(n=k\),
      \(A_{k,k}(\varphi)\) is the average of \(\varphi\) over the \(k!\)
      orderings of \(X_1,\ldots,X_k\).
\item \emph{An identically vanishing average.} If
      \(\varphi:S^2\to\mathbb R\) satisfies
      \(\varphi(y,x)=-\varphi(x,y)\) for all \(x,y\) and
      \(E|\varphi(X_1,X_2)|<\infty\), then \(A_{n,2}(\varphi)=0\) at every
      point of \(\Omega\) and for every \(n\geq2\). Consequently
      \(E(\varphi(X_1,X_2)\mid\mathcal E_n)=0\) almost surely for every
      \(n\geq2\). For real coordinates with \(E|X_1|<\infty\) and
      \(\varphi(x,y)=x-y\), this says
      \(E(X_1\mid\mathcal E_n)=E(X_2\mid\mathcal E_n)\) almost surely.
\item \emph{Repeated indices.} For \(\varphi:S^2\to\mathbb R\) put
      \(\varphi_\Delta(x)=\varphi(x,x)\) and
      \[
       V_n(\varphi)=\frac1{n^2}\sum_{i=1}^n\sum_{j=1}^n\varphi(X_i,X_j).
      \]
      Then, for \(n\geq2\),
      \[
       V_n(\varphi)=\frac{n-1}nA_{n,2}(\varphi)+\frac1nA_{n,1}(\varphi_\Delta).
      \]
      If \(\varphi(X_1,X_2)\) and \(\varphi(X_1,X_1)\) are integrable,
      \[
       V_n(\varphi)=E\Bigl(\frac{n-1}n\varphi(X_1,X_2)
                     +\frac1n\varphi(X_1,X_1)\Bigm|\mathcal E_n\Bigr),
      \]
      and \(V_n(\varphi)-A_{n,2}(\varphi)\to0\) almost surely.
\item \emph{A worked instance.} For i.i.d.\ fair bits,
      \(\varphi(x,y)=xy\) and \(n=2\),
      \(A_{2,2}(\varphi)=X_1X_2=E(X_1X_2\mid\mathcal E_2)\), while
      \(V_2(\varphi)=(X_1+X_2)^2/4\). On the event \(\{X_1+X_2=1\}\),
      of probability \(1/2\), the first is \(0\) and the second is
      \(1/4\). For general \(n\),
      \(A_{n,2}(\varphi)=S_n(S_n-1)/(n(n-1))\) and
      \(V_n(\varphi)=S_n^2/n^2\), where \(S_n=\sum_{j=1}^nX_j\).
\end{enumerate}
\end{hbtheorem}

Item~3 shows exactly what the distinctness of indices in
\(I_{n,k}\) buys. Averaging over all ordered pairs, repetitions
included, gives the conditional expectation of a mixture of
\(\varphi(X_1,X_2)\) with the diagonal value \(\varphi(X_1,X_1)\), with
weight \(1/n\) on the diagonal. The identity of
\cref{thm:p15-exchangeable-averages} at a fixed \(n\) therefore fails for
the repeated-index average, as item~4 shows on an event of probability
\(1/2\), while the two averages have the same almost-sure limit. Durrett
states the average over distinct index sequences and counts them by
\((n)_k\) \cite[p.~252]{DurrettPTEAuthor2019}.

\begin{hbtheorem}{thm:p15-w02c-p15-022-integrability}{proved=annex:p15-w02c-p15-022-integrability}{The integrability hypothesis and the exceptional null sets}
Let \(P\) be exchangeable, \(k\geq1\), \(n\geq k\), and let
\(\varphi:S^k\to\mathbb R\) be measurable.
\begin{enumerate}
\item \emph{The identity holds locally without integrability.} Put
      \[
       W_n=A_{n,k}(|\varphi|),\qquad D_c=\{W_n\leq c\}\quad(c>0).
      \]
      Then \(D_c\in\mathcal E_n\),
      \(E[|\varphi(X_1,\ldots,X_k)|;D_c]\leq c\), the sets \(D_c\)
      increase to \(\Omega\) as \(c\to\infty\), and
      \[
       \mathbf1_{D_c}A_{n,k}(\varphi)
          =E\bigl(\mathbf1_{D_c}\varphi(X_1,\ldots,X_k)\bigm|\mathcal E_n\bigr)
          \quad\text{almost surely}.
      \]
\item \emph{Convergence can survive without integrability.} Let
      \(S=\mathbb R\), let \(\xi\) take the values \(2^m\) and \(-2^m\)
      with probability \(2^{-m-1}\) each, \(m\geq1\), and let \(P\) be
      the law of \((\xi,\xi,\ldots)\). For
      \(\varphi(x_1,\ldots,x_k)=x_1\), the law is exchangeable,
      \(E|\varphi(X_1,\ldots,X_k)|=\infty\), and
      \(A_{n,k}(\varphi)=S_n/n=X_1\) for every \(n\geq k\) almost
      surely. The averages converge almost surely, but neither the averages nor
      their limit is integrable, so the \(L^1\) conclusion and the
      conditional expectations of the theorem have no meaning.
\item \emph{Convergence can fail without integrability.} If the
      \(X_j\) are i.i.d.\ real with \(E|X_1|=\infty\) (for instance with
      the law of \(\xi\) in item~2) and
      \(\varphi(x_1,\ldots,x_k)=x_1\), then
      \(A_{n,k}(\varphi)=S_n/n\), \(P(|X_n|>n\text{ i.o.})=1\), and the
      event that \(A_{n,k}(\varphi)\) has a finite limit has probability
      zero.
\item \emph{The exceptional set depends on the integrand.} Let
      \(S=[0,1]\) with its Borel sigma-field and \(P\) the law of i.i.d.\
      uniform coordinates. For each countable \(D\subseteq[0,1]\),
      almost surely \(A_{n,1}(\mathbf1_D)=0\) for every \(n\) and
      \(E(\mathbf1_D(X_1)\mid\mathcal E)=0\). Nevertheless there is no
      event \(N\) of probability zero such that, for every
      \(\omega\notin N\) and every countable \(D\), the sequence
      \(A_{n,1}(\mathbf1_D)(\omega)\) converges.
\item \emph{Countably many integrands.} For a countable family of
      integrable integrands, one null set serves for all of them.
\end{enumerate}
\end{hbtheorem}

Items~1--3 locate the integrability hypothesis of
\cref{thm:p15-exchangeable-averages}. The averaging identity itself
survives on each invariant event where the averaged absolute values are
bounded; what integrability supplies is a global conditional expectation
and the reverse-martingale convergence. Items~2 and~3 use one and the same
coordinate law: under complete dependence the averages converge
trivially, under independence they have no finite limit. So integrability
is not necessary for almost-sure convergence under every exchangeable
law, but it cannot be dropped from the theorem. The independent case
repeats, for this theorem's averages, the mechanism of item~1 of
\cref{thm:p15-w02b-p15-022-reverse-strong-law-boundary}; the proof is
supplied in full here for the average \(A_{n,k}\).

Items~4 and~5 explain the caveat that different integrands may have
different exceptional null sets. For each fixed countable \(D\) the
theorem holds, but every sample sequence with distinct coordinates is the
witness against its own exceptional set, through a countable set built
from that sequence. Simultaneous statements are therefore made only over
countable families, as in the countable determining class of
\cref{def:p15-directing-law} and the simultaneous recovery of
\cref{thm:p15-directing-law-uniqueness}.

\section{The directing-law space at its edges}
\label{sec:p15-w02c-p15-022-directing}

\begin{hbtheorem}{thm:p15-w02c-p15-022-augmentation}{proved=annex:p15-w02c-p15-022-augmentation}{Null augmentation enlarges the invariant, tail and kernel fields}
Let \(S=\{0,1\}\) with all subsets measurable and let \(P\) be the law of
i.i.d.\ fair bits. Write \(\mathcal N^*=\{A\in\mathcal F:P(A)\in\{0,1\}\}\)
and \(\bar Y=\limsup_nS_n/n\).
\begin{enumerate}
\item \(\mathcal E^{P}=\mathcal T^{P}=\mathcal N^*\).
\item The event \(A\) of \cref{thm:p15-w02c-p15-022-field-strictness}
      with \(B=\{1\}\) satisfies \(A\in\mathcal E\setminus\mathcal T\)
      and \(P(A)=0\). Hence \(A\in\mathcal T^{P}\setminus\mathcal T\),
      and \(\mathcal T^{P}\) is strictly larger than \(\mathcal T\).
\item The constant map and the map
      \[
       K_\omega=\operatorname{Bernoulli}(1/2),\qquad
       K'_\omega=\operatorname{Bernoulli}(\bar Y(\omega))
      \]
      are both regular conditional laws of \(X_1\) given \(\mathcal E\). Their
      fields differ:
      \(\sigma(K)=\{\varnothing,\Omega\}\), while \(\sigma(K')\)
      contains the event \(\{\bar Y=1\}\), which is neither empty nor
      \(\Omega\). Nevertheless
      \(\sigma(K)^{P}=\sigma(K')^{P}=\mathcal N^*\), and
      \(\mathcal E\not\subseteq\sigma(K)\) although
      \(\mathcal E^{P}=\sigma(K)^{P}\).
\end{enumerate}
\end{hbtheorem}

This is the case that the null augmentation in
\cref{def:p15-directing-law} is designed for. The invariant field, the
tail field and the field of a directing law are different collections of
events even for i.i.d.\ fair bits, and the field of the directing law
changes when one version of the regular conditional law is replaced by
another. The equalities of \cref{thm:p15-directing-law-uniqueness} hold
only after every event is allowed to differ from a member of the field by
a null event, and item~3 shows that this is the most that can be asserted.

\begin{hbtheorem}{thm:p15-w02c-p15-022-cocountable}{proved=annex:p15-w02c-p15-022-cocountable}{No countable determining class}
Let \(S=\mathbb R\) and let \(\mathcal A\) be the class of subsets
of \(\mathbb R\) that are countable or have countable complement.
\begin{enumerate}
\item \(\mathcal A\) is a sigma-field containing every singleton.
\item For every countable family \(A_1,A_2,\ldots\) in \(\mathcal A\)
      there are distinct \(x,y\in\mathbb R\) such that, for every \(i\),
      \(x\in A_i\) exactly when \(y\in A_i\). Consequently:
      \begin{enumerate}
      \item no countable family generates \(\mathcal A\);
      \item the point masses \(\delta_x,\delta_y\) are distinct
            probabilities on \(\mathcal A\) that agree on every \(A_i\),
            so no countable subclass of \(\mathcal A\) determines
            probability measures on \(\mathcal A\);
      \item the evaluation sigma-field on the probability measures on
            \((\mathbb R,\mathcal A)\) is not generated by the
            evaluations \(e_{A_i}\) at any countable family.
      \end{enumerate}
\item \((\mathbb R,\mathcal A)\) is not standard Borel.
\end{enumerate}
\end{hbtheorem}

This space is a case in which the standard Borel hypothesis of
\cref{def:p15-directing-law} is needed. The definition chooses a
countable determining pi-system generating \(\mathcal S\), and the
recovery theorem \cref{thm:p15-directing-law-uniqueness} intersects
countably many almost-sure events indexed by it. On
\((\mathbb R,\mathcal A)\) no such class exists, and the evaluation
sigma-field is not reached by countably many evaluations; by
\cref{thm:p15-w02c-p15-022-integrability} item~4, intersecting over an
uncountable family is not available in its place. Whether a regular
conditional law of \(X_1\) given \(\mathcal E\) exists on such a space is
a further question. Durrett states that off nice spaces the
mixture representation fails and refers to Dubins and Freedman (1979) and
Freedman (1980) for counterexamples \cite[p.~254]{DurrettPTEAuthor2019};
the reduction of that failure to the regular conditional law is
\cref{thm:p15-w02b-p15-022-presentable}. Those counterexamples are not
reconstructed here.

\section{The finite-urn bound at its edges}
\label{sec:p15-w02c-p15-022-urn}

\Cref{thm:p15-finite-urn-approximation} bounds, by \(k(k-1)/(2N)\),
the event-supremum total variation distance
\(\|P-Q\|_{\rm TV}=\sup_A|P(A)-Q(A)|\) between the law
\(H_{m,N}^{(k)}\) of the first \(k\) ordered draws without replacement
from \(N\) binary entries with \(m\) ones and the independent law
\(\operatorname{Bernoulli}(m/N)^{\otimes k}\); after mixing over the
final count it gives the same bound for any finite exchangeable binary
vector. Durrett records that the urn laws approach product measure and
that rate bounds can be obtained from this observation
\cite[p.~255]{DurrettPTEAuthor2019}; the computations below are
constructed for this chapter. Write
\[
 d_{m,N}^{(k)}=\left\|H_{m,N}^{(k)}
     -\operatorname{Bernoulli}(m/N)^{\otimes k}\right\|_{\rm TV},
 \qquad
 M_\lambda^{(k)}=\int_{[0,1]}\operatorname{Bernoulli}(\theta)^{\otimes k}
     \,d\lambda(\theta)
\]
for \(0\leq m\leq N\), \(1\leq k\leq N\) and a probability \(\lambda\)
on \([0,1]\). The inputs are the population \((m,N)\), the draw count
\(k\) and, for the mixed bound, the law of a finite exchangeable binary
vector \((X_1,\ldots,X_N)\) with \(V=S_N/N\) and \(\lambda_N\) the law
of \(V\). The conclusions are exact values of the distances, compared
with the bound.

\begin{hbtheorem}{thm:p15-w02c-p15-022-urn-boundary}{proved=annex:p15-w02c-p15-022-urn-boundary}{The finite-urn bound computed, vacuous and unattainable}
Let \(p=m/N\).
\begin{enumerate}
\item \emph{Two draws.} For \(N\geq2\),
      \[
       d_{m,N}^{(2)}=\frac{2p(1-p)}{N-1}=\frac{2m(N-m)}{N^2(N-1)}
       \leq\frac{1}{2(N-1)},
      \]
      with the last inequality an equality exactly when \(N=2m\). The
      distance equals the bound \(1/N\) exactly when \(N=2\) and
      \(m=1\); for every \(N\geq3\) it is strictly smaller, and it is at
      most \(N/(2(N-1))\) times the bound.
\item \emph{Computed instances.} For \((m,N)=(5,10)\),
      \(d_{5,10}^{(2)}=1/18\) against the bound \(1/10\). For
      \((m,N)=(2,4)\) and \(k=2,3,4\),
      \[
       d_{2,4}^{(2)}=\tfrac16,\quad
       d_{2,4}^{(3)}=\tfrac14,\quad
       d_{2,4}^{(4)}=\tfrac58,
       \qquad\text{against the bounds}\qquad
       \tfrac14,\quad\tfrac34,\quad\tfrac32 .
      \]
      For fixed \((m,N)\), \(d_{m,N}^{(k)}\) is nondecreasing in \(k\).
\item \emph{All draws.} For \(k=N\),
      \[
       d_{m,N}^{(N)}=1-\binom Nm p^m(1-p)^{N-m}.
      \]
      If \(0<m<N\) this lies strictly between \(0\) and \(1\); if
      \(m\in\{0,N\}\) then \(d_{m,N}^{(k)}=0\) for every \(k\). For
      \((m,N)=(5,10)\), \(d_{5,10}^{(10)}=193/256\) against the bound
      \(9/2\).
\item \emph{The vacuous range.} The bound is at least \(1\) exactly
      when \(k(k-1)\geq2N\); every event-supremum distance between
      probabilities is at most \(1\), so in that range the theorem
      asserts nothing. For \(N=4\) this is \(k=4\) only, and for
      \(N=10\) it is \(5\leq k\leq10\).
\item \emph{The mixed bound at two draws.} For a finite exchangeable
      binary vector with \(N\geq2\),
      \[
       \left\|\mathcal L(X_1,X_2)-M_{\lambda_N}^{(2)}\right\|_{\rm TV}
       =\frac{2E[V(1-V)]}{N-1}.
      \]
      \begin{enumerate}
      \item If \(N=2\) and \(X_1,X_2\) are independent fair bits, the
            distance is \(1/4\), although
            \(\mathcal L(X_1,X_2)=M_{\delta_{1/2}}^{(2)}\) is itself a
            binary mixture.
      \item If \(N\) is even and \((X_1,\ldots,X_N)\) is uniform over the
            words with exactly \(N/2\) ones, then for every probability
            \(\lambda\) on \([0,1]\),
            \[
             \left\|\mathcal L(X_1,X_2)-M_{\lambda}^{(2)}\right\|_{\rm TV}
             \geq\frac{1}{2(N-1)},
            \]
            with equality at \(\lambda=\delta_{1/2}=\lambda_N\). For
            \(N=2\) the minimum \(1/2\) equals the bound
            \(k(k-1)/(2N)\).
      \end{enumerate}
\end{enumerate}
\end{hbtheorem}

Item~1 shows that the order \(1/N\) of the bound is the true order at
\(k=2\): for fixed \(p\) strictly between \(0\) and \(1\),
\(N\,d_{m,N}^{(2)}\) tends to \(2p(1-p)\) along populations with
\(m/N=p\), while the constant in the bound is \(1\). The bound is
attained only at the smallest population, and it ignores \(m\), whereas
the true distance vanishes at \(m\in\{0,N\}\) for every \(k\) (item~3).
Items~3 and~4 show where the bound carries no information: once
\(k(k-1)\geq2N\), as at \(k=N\geq3\), the true distance is still
strictly below \(1\) when \(0<m<N\) and is given by the binomial formula,
but the theorem does not reach it.

Item~5 bounds the mixed statement. Its mixing law is the law of the
empirical frequency \(S_N/N\), not a best mixing law, so the
approximation error can be positive for a vector that is exactly an
independent mixture (item~5(a)). Some finite exchangeable vectors are at
positive distance from every independent mixture (item~5(b)), so no
choice of mixing law makes the mixed bound exact, and at \(N=k=2\) the
bound \(1/2\) cannot be lowered for that vector even by choosing the
mixing law freely. The vector of item~5(b) with \(N=2\) is the
two-observation law of \cref{thm:p15-carnap-history-finite} that has no
exchangeable extension to three observations; item~5(b) measures how
far it is from the independent mixtures.

\section{Relationships and sources}
\label{sec:p15-w02c-p15-022-relationships}

\Cref{thm:p15-w02c-p15-022-field-strictness,thm:p15-w02c-p15-022-not-exchangeable}
bound \cref{def:p15-exchangeable-canonical} and
\cref{thm:p15-exchangeable-canonical}: the first decides strictness of the
two inclusions literally and modulo null events and shows that
\(\mathcal E^{P}=\mathcal T^{P}\) needs exchangeability; the second gives
a law with identical coordinate laws and a symmetric initial pair that is
not exchangeable, and shows what the canonical reduction does not
transfer. \Cref{thm:p15-w02c-p15-022-average-edges} bounds the averages
\(A_{n,k}\) of the definition, and
\cref{thm:p15-w02c-p15-022-integrability} bounds
\cref{thm:p15-exchangeable-averages} in its integrability hypothesis and
its per-integrand exceptional sets.
\Cref{thm:p15-w02c-p15-022-augmentation,thm:p15-w02c-p15-022-cocountable}
bound \cref{def:p15-directing-law}: the first gives null augmentations
strictly larger than the fields they augment, and the second a value space
on which the countable determining class required there does not exist.
\Cref{thm:p15-w02c-p15-022-urn-boundary} bounds
\cref{thm:p15-finite-urn-approximation}: it computes the distance the
theorem bounds, locates the populations at which the bound is attained
and the draw counts at which it is vacuous, and shows that the mixed bound
for finite exchangeable vectors cannot in general be made exact.

The Hewitt--Savage law used in
\cref{thm:p15-w02c-p15-022-augmentation} is
\cref{thm:p15-hewitt-savage-zero-one}, and the second Borel--Cantelli
lemma used in \cref{thm:p15-w02c-p15-022-integrability} is
\cref{thm:p15-borel-cantelli-second}. Durrett's definitions of
\(\mathcal E_n\) and \(\mathcal E\) are on p.~250, the
distinct-index average with its count \((n)_k\) in Lemma~4.7.7 on
p.~252, and the observation that \(\mathcal E\) may be nontrivial for an
exchangeable sequence, so that the limit of the averages need not be a
constant, on p.~253 \cite[pp.~250--253]{DurrettPTEAuthor2019}. Durrett
states Theorem~4.7.10, the binary case of the representation, on
p.~254, and his remark that the urn laws approach product measure, with
bounds on the rate, is on p.~255 \cite[pp.~254--255]{DurrettPTEAuthor2019}. The
examples and their proofs here are constructed for this chapter; they
are not presented as Durrett's.
